%! TEX root = penilaiandiri.tex
% the main TeX file which is intended to compile, :VimtexReload after adjustment
% :h vimtex-tex-root
\documentclass[../main.tex]{subfiles}
%----------------------------------------------------------------------------------------
%	TITLE SECTION
%----------------------------------------------------------------------------------------
% MAIN TITLE SECTION
\title{
\textbf{Penilaian diri} \\
} % Title and subtitle
%\date{\textbf{\DTMtoday}}
\date{\textbf{\today}}
\author{uliah}
%\begin{tabular}{@{}ll@{}}
%	Nama & : Muhammad Uliah Shafar\\
%	NIM & : 21020119420029\\
%\end{tabular}
%}

%----------------------------------------------------------------------------------------
% OTHER TITLE SECTION

%\title{\textbf{Sistem Sarana dan Prasarana Jl. Pinggir Laut} \\ {\Large\itshape Infrastructure of Waterfront Parepare City}} % Title and subtitle

%\author{\textbf{Uliah Shafar} \\ \textit{Universitas Diponegoro}} % Author and institution

%\date{\today} % Date, use \date{} for no date

%----------------------------------------------------------------------------------------



\begin{document}

%\maketitle % Print the title section

%----------------------------------------------------------------------------------------
%	ABSTRACT AND KEYWORDS
%----------------------------------------------------------------------------------------
%\begin{abstract}
%Lorem ipsum dolor sit amet, consectetur adipiscing elit.
%\end{abstract}

%\hspace*{3.6mm}\textit{Keywords:} Lorem, Ipsum % Keywords

\vspace{30pt} % Vertical whitespace between the abstract and first section

%----------------------------------------------------------------------------------------
%	ESSAY BODY
%----------------------------------------------------------------------------------------


Saya sangat semangat untuk belajar hal baru terutama yang menurut saya akan memudahkan dan mempercepat pekerjaan saya di masa depan. Hal baru yang terakhir saya pelajari adalah latex, bahasa markup untuk penulisan karya ilmiah.

Separuh mahasiswa arsitektur saya temukan tidak dapat melanjutkan studi arsitektur di kampus saya. Meskipun saya bukan mahasiswa dengan nilai tinggi sewaktu sarjana, saya tetap menekuni arsitektur hingga lulus dengan kegigihan.

Saya merasa memiliki etos kerja. Saya selalu memiliki motivasi untuk melakukan yang terbaik hingga tenggat waktu. Meskipun hasil tidak selalu sesuai dengan harapan, keringat dan air mata merupakan penghargaan terbesar bagi saya.

Saya memiliki kemampuan bahasa inggris yang memadai. Hasil tes bahasa inggris ielts saya yang terakhir adalah C1 atau peringkat menengah. Tesis saya baru-baru ini memiliki sitasi jurnal bahasa inggris diatas 50\%

Saya telah menekuni satu bidang yakni arsitektur selama hampir satu dekade. Dalam kurung waktu tersebut saya menempuh pendidikan arsitektur dari SMK hingga perguruan tinggi secara linear. Dan juga pernah magang dan bekerja di biro arsitektur.


Kelemahan yang dominan dalam diri saya adalah sifat terburu-buru. Saya ingin menuntaskan segala hal yang ada di pikiran saya dalam satu waktu. Namun, itu kadang tertangani dengan menetapkan prioritas.

Saya pikir saya termasuk orang yang goal-oriented. Ini menyebabkan saya ingin menyelesaikan gagasan secepat mungkin tanpa mendengar pemikiran orang lain. Sehingga, saya mulai belajar tenang dan selalu memikirkan konsekuensi dari tindakan.

Saya biasanya kurang detail dalam mengerjakan sesuatu. Ini mungkin karena target yang banyak. Maka saya memilih satu tugas yang penting untuk dikerjakan dengan maksimal.

Kadang dalam pekerjaan, saya menjadi perfectionism. Saya terlena dengan satu pekerjaan hingga lupa tenggat waktu. Saya mengatasinya dengan belajar menemukan keseimbangan antara kesempurnaan, sangat bagus dan tepat waktu.

\subsection{Pengalaman yang membanggakan}
Pengalaman yang membanggakan bagi saya baru-baru ini adalah ketika saya berhasil mendapatkan IPK 4.00 di magister arsitektur Universitas Diponegoro.  Selama mengejar gelar magister, saya mengawali pekerjaan di sebuah kantor kontraktor yang saat itu merenovasi sekolah yang merupakan kebanggan juga.
Pengalaman yang membanggakan selain itu adalah ketika saya ikut serta menyukseskan kegiatan-kegiatan organisasi pada saat kuliah seperti sukarelawan atau kompetisi. Beberapa pengalaman lainya seperti pernah menjadi asisten dosen, mendapatkan nilai IELTS 5.5, dan menjadi bagian tim basket merupakan kebanggan tersendiri untuk saya.

Pengalaman yang membanggakan yang sedikit tergerus oleh waktu adalah ketika saya mengikuti beragam kegitan organisasi di SMK. Saya pernah ikut dalam kegiatan LDK (latihan dasar kepemimpinan) oleh tentara dan kelopak pramuka. Pada saat kelopak saya menjuarai lomba film pendek se-Sulawesi Selatan. Diluar kegiatan organisasi, saya sempat magang di perusahaan semen tonasa.

Pengalaman yang terakhir tapi tidak kalah penting adalah ketika saya diterima mejadi dosen di perguruan tinggi Universitas Muhammadiyah Parepare sebagai dosen Perencanaan Wilayah dan Kota.

\subsection{Pengalaman yang kurang membanggakan}

Pengalaman yang kurang membanggakan adalah ketika saya terlambat mengumpulkan tugas proposal tugas akhir. Itu membuat saya harus mengulang satu tahap (setengah semester) dan mempengaruhi nilai saya. Pengalaman ini merupakan bagian dari kelemahan saya yakni perfectionism. Saya berusaha membuat tugas saya sangat sempurna hingga mengabaikan tenggak waktunya. Kedepannya saya akan menyeimbangkan antara kesempurnaan, sangat bagus dan tepat waktu. Selain itu, saya akan menempatkan hal-hal penting lebih dulu daripada yang tidak penting.


\subsection{Apa hal terakhir yang anda ajarkan pada diri anda sendiri?}

Terus berusaha, tidak perlu khawatir apakah kita akan mendapatkan atau tidak, apakah jalannya sudah benar atau tidak. Apapun yang ditakdirkan, kita pasti dapatkan begitupun sebaliknya. Saya ingat perkataan ahli psikologi, Jordan Peterson, ``Teruslah berjalan, jika ditengah jalan kita berpikir itu salah, kita sedang berproses. Kita tidak akan tahu itu salah apabila kita tidak pernah berjalan setengah jalan yang salah."


\subsection{Apakah anda pernah melakukan tugas di luar ruang lingkup tugas anda yang seharusnya? jelaskan}
Dunia memang selalu menuntut untuk fleksibel. Oleh karena itu saya selalu berusaha untuk beradaptasi.
Salah satu contoh konkret belum lama ini terjadi pada saat saya bekerja.
Saya mengoperasikan drone dan menyajikan foto hasil drone tersebut.
Tujuan tugas tersebut adalah membuat laporan tentang foto nol dan proses hingga selesai bangunan yang sedang dikerjakan.
Sama halnya terjadi pada saat saya mengerjakan tesis.
Saya membuat video dari objek penelitian yang saya teliti.
Tujuan tugas tersebut adalah untuk mengenalkan lebih dalam profesor saya terkait lokasi penelitian tesis saya.


\subsection{Kesalahan apa yang pernah anda lakukan selama belajar/bekerja?}
Kesalahan di masa lalu bukanlah vonis seumur hidup, tetapi itu adalah pembelajaran yang berharga.
Kesalahan saya sewaktu belajar adalah menargetkan banyak tujuan diluar kemampuan saya.
Tujuan-tujuan tersebut adalah mengambil mata kuliah yang penuh, aktif dalam berorganisasi, mengejar kesempatan untuk lanjut kuliah, dan mendaftar sayembara desain arsitektur.
Tujuan tersebut saya laksanakan dalam satu waktu tanpa memikirkan yang mana perlu diutamakan. Alhasil, saya kewalahan dan tidak ada hasil yang signifikan. Pembelajaran yang saya ambil adalah saya harus menargetkan satu tujuan dalam satu waktu dalam bahasa inggrisnya \textit{one at the time}. Saya harus memilih satu hal yang perlu saya fokuskan dengan maksimal. Sisanya saya kesampingkan untuk sisa tenaga yang saya miliki diakhir hari.

\subsection{Sebutkan dan jelaskan 1 (satu) hal yang membedakan anda dengan peserta lainnya?}

Saya berasal dari kota Parepare, salah satu kota diantara dua kota di Sulawesi Selatan. Kota yang merupakan bagian dari wilayah Republik Indonesia yang memerlukan pengembangan, meskipun bukan termasuk daerah afirmasi. Tempat dimana saya kemungkinan akan mengabdi dan mencurahkan segenap manfaat yang bisa saya lakukan. Semuanya itu saya lakukan untuk meningkatkan daerah kelahiran saya, tempat keliharan semua yang saya cintai, yang saya sangat cintai dan banggakan.



%Contoh lainnya adalah desain dengan prinsip energi adalah membangun ruang publik dengan konsumsi energi rendah menggunakan panel surya.


%Berdasarkan konsep keberlanjutan, desain dapat menjawab beberapa tantangan dari urbanisasi. Pertama, tantangan konsumsi sumber daya yang meningkat. Desain dapat menggunakan panel surya dan material yang dapat mengurangi konsumsi energi. Sehingga tercipta \textit{zero energy}. Kedua, tantangan kualitas lingkungan yang menurun.


%Banyak arsitek saat ini menyerukan usaha \textit{zero energy} pada bangunan. Usaha tersebut dapat terealisasikan di ruang publik melalui penggunaan panel surya. Selain itu,

%Itu dapat diimplementasikan pada ruang publik melalui penggunaan panel surya. Lebih lanjut,

%Implementasinya pada ruang publik yaitu penggunaan panel surya.

%Arsitek harus pandai dalam menjawab tantangan urbanisasi.
%Perancang dapat
%Berdasarkan konsep keberlanjutan, perancang dapat menjawab sejumlah tantangan.
%
%Ber
%
%Ta

%Beberapa diantaranya adalah z
%Mereka dapat mendesain sebuah ruang publik dengan menggunakan tekologi yang muthahir

\begin{comment}
%Desain ruang publik
Sejumlah kota telah melakukan desain ruang publiknya. Salah satunya adalah jepara.

Beberapa kota telah mendesain ruang publik mereka dengan tujuan meningkatkan keberlanjutan kota. Salah satunya adalah jepara. Jepara menjadi kota di jawa yang menjadi pioner dalam segi keberlanjutan. Jepara telah melakukan sejumlah langkah dalam proses keberlanjutan. Satu yang terpenting adalah mendesain ruang publik seperti tepi laut, pantai dalam kepariwisataan.

What you have got to design public space for sustainability city

Saya dapat menjadi bagian dari gerakan tersebut dengan menempuh pendidikan yang lebih tinggi.
Melalui pendidikan lebih tinggi, saya akan mempelajari lebih dalam konsep keberlanjutan dalam sebuah ruang publik.
Adapaun pendidikan yang saya tempuh adalah tingkat doktor di univeristas dalam negeri. Dari tingkat sarjana hingga magister, saya telah mempelajari konsep keberlanjutan dalam ilmu arsitetur, harapan saya untuk studi ke jenjang doktor adalah menggali lebih dalam ilmu perancangan sebuah kota khususnya ruang publik. Saat ini, pendidikan arsitektur di universitas-universitas dalam negeri sudah snagat baik.


%??? benar merah urbanisasi dengan konsep keberlanjutan dalam mendesain ruang publik.


\end{comment}

\begin{comment}

Suatu kebanggaan apabila saya dapat menyebarkan ilmu pengetahuan yang saya pelajari.
Pembelajaran terkait arsitektur masih cukup jarang di perguruan tinggi daerahku dan sejumlah daerah tetangga.
Satu yang dapat mengakomodasi pembelajaran tersebut adalah meningkatnya kompetensi dosen.

Kebanyakan program studi arsitektur berpusat di kota besar Makassar. Saya berpendapat penyebaran ilmu pengetahuan ke sejumlah kota akan memberikan wawasan bagi mereka yang tidak berkesempatan berkuliah di tempat yang terlampau jauh dari kampung mereka.

kepada lulusan PT dari beragam daerah mengenai bagaimana perancangan kota semestinya. Itu akan


%Meskipun demikian, beberapa perguruan tinggi baru-baru ini membuka program studi terkait arsitektur, seperti tempat saya mengajar.



\rule{\linewidth}{0.5pt}




Hiruk pikuk, getar getir dan siang malam dunia pendidikan arsitektur alhamdulillah (puji tuhan) telah saya lalui.

Ketika saya merantau ke pulau Jawa untuk kuliah S1, tujuh tahun yang lalu, saya mempelajari banyak hal-hal yang baru. Beberapa diantaranya adalah pengetahuan, tempat, orang-orang dan suasana. Secara tidak langsung, itu menambah pengalaman dan rasa ingin tahu saya. Apapun itu membuat saya bertekad untuk kembali melanjutkan studi. Saat ini saya telah menempuh gelar sarjana di Universitas Teknologi Yogyakarta dan magister di Universitas Diponegoro dengan jurusan Arsitektur.

Dengan kualifikasi S1 Arsitektur di Yogyakarta dan S2 Arsitektur Undip di Semarang, saya percaya masih mampu meraih pengalaman yang mungkin masih berupa angan-angan tersebut dalam bentuk studi S3.

Alasan lain yang membulatkan tekad itu adalah umur yang relatif muda. Seperti pepatah yang mengatakan belajar di waktu kecil bagai mengukir diatas batu, sedangkan sesudah besar bagaiakan melukis diatas air. Banyak penelitian juga mengindikasikan potensi penurunan kognisi berkaitan dengan peningkatan usia. Dengan begitu, semoga studi yang nantinya saya lalui dapat berjalan dengan tempo yang tepat.
Selain itu, kota yang saya tinggali, Parepare, merupakan kota yang sedang berkembang pesat. Daerah yang saya tempati yang agak jauh dari pusat kota ini mulai terbangun banyak perumahan. Hal itu berdampak pada peningkatan bisnis, ekonomi, dan pariwisata di pusat kota dan sekitarnya. Ketika saya masih duduku di bangku SMK, di SMK negeri 2 Parepare, saya biasanya berjalan-jalan baik bersama keluarga atau teman sekolah. Waktu itu, kota hanya memiliki sedikit pusat keramaian. Ornamen jalan pun masih sedikit dan belum ada lampu-lampu yang gemerlap.
Berbeda dengan saat ini, dimana suasana kota mulai mengikuti perkembangan kota metropolitan. Dimana-mana terdapat lampu-lampu, lanmark, dan tempat rekreasi, seperti monumen BJ Habibie, Tepi Laut Tonrangeng, dan Taman Andi Makkasau. Kadang ketika saya jalan bersama keluarga, ketiga adik dan kedua orangtua, kami terpukau melihat gemerlap lampu-lampu dijalan.



%%%

Selain itu, kota yang saya tinggali ini merupakan kota yang sedang ``rising".

%Penelitian terkait umur dan kognisi belakangan ini muncul yang mana mengindikasikan adanya korelasi terhadap peningkatan umur dan penurunan kognisi.

Pendidikan adalah senjata yang paling ampuh.
Otak saya dapat menerima

%Terlahir di Parepare, sebuah kota kecil di pulau yang berbentuk huruf k, lokasi geografisnya tepat di utara kota daeng sekitar 150km jauhnya. Kota daeng atau Makassar adalah ibukota dari provinsi Sulawesi Selatan.

%This file powered by \citep{einstein} and \citep{latexcompanion}. It can also be found in \cite{knuthwebsite}.

\end{comment}

%----------------------------------------------------------------------------------------
%	BIBLIOGRAPHY
%----------------------------------------------------------------------------------------

%\bibliographystyle{apalike}

%\bibliography{biblio.bib}

\end{document}
